\documentclass{article}
\usepackage{amsmath,hyperref}
\title{Joint coverage and radio recovery in an ordered corridor}
\date{}
\begin{document}
\maketitle

\paragraph{Physical and service contract.}
There are $n$ ordered robots on a line with known initial positions
$p_i$, maximum speeds $v_i$ and no acceleration constraint. Each
robot must stay within an assigned closed interval, already reduced
for fixed obstacle clearance. Consecutive bodies must keep separation
at least $g_i>0$ throughout. They start in that safe set.
At deadline $T$ their positions must provide connected radio and cover
the closed service interval $[L,U]$, despite any set of at most $k$
further deletions. A deletion removes radio and sensing, not initial
clearance or the physical occupied position. The calculation is about
a stationary test at $T$, not continued tolerance under new motions.
All robots have the same symmetric radio range $R$ and the same
interval sensing radius $b$. Require $n\ge k+2$. No wireless
propagation, detection error or ground-truth position inference is
claimed.

\paragraph{Exact placement.}
Intersect each assigned interval with $[p_i-v_iT,p_i+v_iT]$.
For a point in the resulting box, the inequalities
$x_{i+1}-x_i\ge g_i$ are necessary and sufficient for safe continuous
motion to it by $T$. For sufficiency interpolate linearly between
the initial and terminal configurations; both the box of hard
position limits and the separation constraints are convex. If $T=0$,
the reachable box already fixes the initial configuration.
No discrete time approximation is used.

Write $q=k+1$. Robust radio connectivity holds exactly when
$x_{i+q}-x_i\le R$ for all $1\le i\le n-q$. Necessity deletes the
$k$ intervening robots: no other pair crosses that gap with a shorter
distance. Sufficiency holds because two consecutive survivors span
at most $q$ original steps and can be included in a $q$-step window.
Robust pointwise sensing holds exactly when every point in $[L,U]$
is seen by at least $q$ robots before deletion. The blocks of
$q$ geometrically consecutive robot positions suffice for that test.
Their common coverage intervals are
$J_i=[x_{i+q-1}-b,x_i+b]$; an inverted interval contributes nothing.
Both endpoints are nondecreasing in $i$. The union covers $[L,U]$
exactly when there exist $i\le j$ with $L\in J_i$, $U\in J_j$ and
$x_{t+q}-x_t\le 2b$ for $i\le t<j$. A gap between successive
non-overlapping intervals cannot be repaired by skipping an interval
with both endpoints between them. If one interval covers both ends,
take $i=j$. This includes singleton service zones.

For each of the $O(n^2)$ choices of $i,j$, all of the conditions are
ordinary difference constraints, with $O(n)$ terms. Shortest paths
give potentials satisfying every condition, or a negative cycle
proves infeasibility in that case. The union test is exact over the
reals. Integral inputs yield an integral feasible witness, if
one exists, so exhaustive integer placement is a valid independent
check for those fixtures. The solver permits the same loss set in
both requirements. Since both must hold for every set, neither may
borrow coverage or radio from a different placement.

\paragraph{An incompatible pair of duties.}
Three robots must tolerate one loss. One robot is already assigned to
position 0, the others may reach coordinates in $[1,9]$, with unit
clearance. Let $R=4$, $b=2$ and $[L,U]=[7,8]$. The radio-only task
can use $(0,1,2)$; coverage alone can use $(0,6,7)$.
Together they demand $x_3\le4$ and $x_3\ge6$.
Conversely a successful terminal solution does not imply that sensing
can be maintained during interpolation: two robots at $(0,10)$
and later $(-10,0)$ each cover the origin with $b=1$, but the straight
midpoint $(-5,5)$ does not. This is why the coverage obligation in
the model is explicitly deadline recovery.

\paragraph{An information cutoff, not a physical impossibility.}
For $q\ge2$, take $2q-1$ robots initially centered at equally spaced
points with gap $g$, all with the same speed $v>0$. Assume the corridor
permits the translations below, $R\ge qg$, and $(q-1)g\le2b$.
The protected point is either $-(b+d)$ or $+(b+d)$ for $d>0$;
all receive a reliable common indication at time $\tau$. Before it,
their histories are identical and there is no informative action.
One branch requires the median robot at or left of $-d$ at $T$,
the other at or right of $d$, since each requires a strict majority
within sensing distance. From any common median at $\tau$ at
least one branch requires travel of at least $d$. Thus
$v(T-\tau)\ge d$ is necessary. It is sufficient to wait at the
initial formation and translate by $-d$ or $+d$ after the indication.
The appropriate $q$ sensors then cover the protected point; radio
resilience and clearances hold during the entire translation.
Both worlds can be independently feasible even when their common
information arrives after this tight cutoff. The proof does not extend
to private, delayed or disagreeing indications without a new model.

\paragraph{Comparison.}
This adds sensing coverage to the existing ordered
\href{https://github.com/eternal-november/eos-research-programmes/pull/375}{radio-resilience test}.
The separated survivor-set oracle checks continuous interval coverage,
not only sampled points. The constructive rule matches it and applies
without enumerating deletion sets. It uses standard difference
constraints, not a gain over a similarly factored direct planner.
The physical semantics are far narrower than heterogeneous
\href{https://arxiv.org/abs/2210.01354}{barrier-function synthesis}
or geometric
\href{https://www.cs.memphis.edu/~santosh/Papers/BarrierCoverageMobiCom05.pdf}{barrier coverage}.
This is a rigorous restricted R6 comparison, not a field result
or the independent enclosed-space mission.
\end{document}
